%%%%%%%%%%%%%%%%%%%%%%%%%%%%%%%%%%%%%%%%%%%%%%%%%%%%%%%%%%%%%%
%%%%%%%%%% MA-0781 Pasantía en matemáticas %%%%%%%%%%%%%%%%%%%
%%%%%%%%%%%%%%%%%%%%%%%%%%%%%%%%%%%%%%%%%%%%%%%%%%%%%%%%%%%%%%
\newpage
\subsection*{MA-0781 Pasantía en matemáticas}
\addcontentsline{toc}{subsection}{MA-0781 Pasantía en matemáticas}

\hypertarget{MA0781}{}
\begin{refsection}
\begin{table}[H]
\centering{
\begin{tabular}{|ll|}
\hline
\textbf{Año:} IV  & \textbf{Requisitos:} Haber aprobado al menos uno de: MA-0661 \\ & Algebra Abstracta II, MA-0635 Introducción a la Topología \\ & o MA-0641 Análisis Numérico I \\ 
\textbf{Ciclo:} VIII  & \textbf{Correquisitos:} Ninguno \\ 
\textbf{Tipo de curso:} Práctico & \textbf{Horas presenciales por semana:} Según convenio \\ 
\textbf{Créditos:} 4 & \textbf{Horas de trabajo independiente por semana:} Según convenio \\ \hline
\end{tabular}
}
\end{table}

\subsubsection*{Descripción del curso}

Este curso optativo ofrece a la persona estudiante un espacio de formación práctica en el que aplica conocimientos matemáticos en un entorno real de trabajo, bajo supervisión académica del Departamento de Matemáticas. El curso admite \textbf{dos modalidades}, que se eligen al matricular y se formalizan mediante un plan de trabajo aprobado por la coordinación del curso:

\begin{enumerate}
\item \textbf{Pasantía académica:} la persona estudiante realiza trabajos de investigación elemental---por ejemplo, revisión documental, experimentación numérica, programación científica, elaboración de notas o informes técnicos---bajo la dirección de una persona docente o investigadora de la Universidad de Costa Rica o de otra universidad o instituto de investigación. El énfasis está en el acercamiento a la práctica investigativa y en el desarrollo de autonomía intelectual.

\item \textbf{Pasantía profesional:} la persona estudiante realiza un trabajo en un ente externo al ámbito académico (empresa, institución estatal, organización no gubernamental u otro organismo con necesidades cuantitativas o de modelación). El énfasis está en la aplicación de herramientas matemáticas a problemas del entorno laboral y en la comunicación con equipos no necesariamente especializados en matemática.
\end{enumerate}

En ambas modalidades, la persona estudiante cuenta con una \textbf{persona supervisora en el sitio de la pasantía} (docente, investigadora o profesional del ente receptor) y con una \textbf{persona docente coordinadora} del Departamento de Matemáticas, responsable de la matrícula, del aval del plan de trabajo y de la evaluación final del curso.

\subsubsection*{Objetivos}

\textbf{General:} Que la persona estudiante integre y aplique conocimientos matemáticos en un entorno académico o profesional, desarrollando autonomía, responsabilidad y capacidad de comunicar resultados.

\textbf{Específicos:} Al finalizar el curso se espera que la persona estudiante sea capaz de:

\begin{enumerate}
\item Formular, de acuerdo con la modalidad elegida, un plan de trabajo acotado, realista y alineado con su formación matemática.
\item Ejecutar tareas de investigación elemental o de aplicación profesional con rigor, ética y cumplimiento de plazos.
\item Utilizar herramientas matemáticas, computacionales o documentales pertinentes al problema abordado.
\item Comunicar oralmente y por escrito los resultados, limitaciones y aprendizajes de la pasantía ante la coordinación del curso y, cuando corresponda, ante el ente receptor.
\item Reflexionar sobre el papel de la matemática en contextos de investigación o de ejercicio profesional.
\end{enumerate}

\subsubsection*{Contenidos}

Los contenidos \textbf{no se fijan de manera permanente} en el plan de estudios: dependen de la modalidad y del plan de trabajo aprobado para cada persona estudiante. Al inicio del ciclo lectivo (o antes de la matrícula, según indique la coordinación), deben quedar documentados por escrito:

\begin{itemize}
\item la modalidad (académica o profesional);
\item el sitio de la pasantía y la persona supervisora en el sitio;
\item los objetivos específicos del plan de trabajo;
\item las actividades previstas y el cronograma tentativo;
\item los productos esperados (informe, código, notas, presentaciones u otros);
\item los criterios de seguimiento y evaluación.
\end{itemize}

A título orientativo, las actividades pueden incluir:

\begin{enumerate}
\item \textbf{En la pasantía académica:}
\begin{enumerate}
    \item Revisión bibliográfica y síntesis documental sobre un tema acotado.
    \item Implementación o experimentación numérica y programación científica.
    \item Elaboración de notas, reportes técnicos o borradores de comunicación científica.
    \item Participación en reuniones de grupo de investigación o seminarios del grupo anfitrión.
\end{enumerate}

\item \textbf{En la pasantía profesional:}
\begin{enumerate}
    \item Análisis, modelación o procesamiento de datos en el contexto del ente receptor.
    \item Apoyo en la formulación o validación de modelos cuantitativos.
    \item Documentación técnica y comunicación de resultados a equipos interdisciplinarios.
    \item Cumplimiento de normas de confidencialidad y ética profesional del ente receptor.
\end{enumerate}
\end{enumerate}

\subsubsection*{Metodología}

El curso se desarrolla principalmente \textbf{fuera del aula}, en el sitio de la pasantía, con acompañamiento de la persona supervisora en el sitio y supervisión académica de la persona docente coordinadora.

\noindent \textbf{Lineamientos metodológicos:}

\begin{enumerate}
    \item Antes de iniciar las actividades, la persona estudiante debe contar con un \textbf{plan de trabajo} avalado por la coordinación del curso y aceptado por la persona supervisora en el sitio.
    \item La pasantía académica requiere que la persona supervisora en el sitio sea docente o investigadora de la UCR o de otra universidad o instituto de investigación, con formación o trayectoria pertinente al tema.
    \item La pasantía profesional requiere un convenio, carta de aceptación o mecanismo institucional equivalente entre el Departamento de Matemáticas (o la unidad académica competente) y el ente receptor, cuando así lo exija la normativa vigente.
    \item La persona estudiante mantiene un registro de actividades (bitácora o equivalente) y reporta avances a la coordinación según el calendario acordado.
    \item Al finalizar, la persona estudiante entrega un \textbf{informe final} y realiza una \textbf{presentación oral} ante la coordinación del curso (y, si se estima conveniente, ante la persona supervisora en el sitio).
    \item La persona supervisora en el sitio emite una evaluación escrita del desempeño, que la coordinación incorpora en la calificación del curso.
\end{enumerate}

\textbf{Sugerencias metodológicas:}

\begin{enumerate}
    \item Acotar el alcance del plan de trabajo al tiempo disponible en el ciclo lectivo, privilegiando un producto concreto y evaluable.
    \item En la modalidad académica, priorizar problemas que permitan un primer contacto con la investigación (lectura dirigida, experimentos reproducibles, documentación clara) sin pretender resultados de publicación obligatoria.
    \item En la modalidad profesional, explicitar desde el inicio las restricciones de confidencialidad y los formatos de entrega aceptables para el informe académico.
    \item Programar al menos una reunión intermedia de seguimiento entre la persona estudiante, la persona supervisora en el sitio y la coordinación del curso.
\end{enumerate}

\subsubsection*{Evaluación}

La evaluación es responsabilidad de la persona docente coordinadora del curso, con base en el plan de trabajo aprobado y en los insumos de la persona supervisora en el sitio. Se sugiere la siguiente distribución, ajustable según la naturaleza de cada pasantía:

\begin{enumerate}
\item Plan de trabajo y cumplimiento de hitos intermedios: 20\,\%.
\item Evaluación de la persona supervisora en el sitio: 30\,\%.
\item Informe final escrito: 30\,\%.
\item Presentación oral final: 20\,\%.
\end{enumerate}

La no entrega del informe final o la ausencia injustificada a la presentación oral implican la no aprobación del curso, salvo lo dispuesto por el Reglamento del Régimen Académico Estudiantil y demás normativa institucional vigente.

\subsubsection*{Bibliografía}

No aplica de manera fija. Las referencias se definen en el plan de trabajo de cada pasantía, según el tema y la modalidad.

\end{refsection}
